\subsubsection{Topological models used by the evaluator}
\label{subsec:topological-evaluator-models}

\begin{definition}[Cut-open tangle at a dotted component]
\label{def:cut-open-dotted-tangle}
Let \(\mathcal H\) be a one-ended Kirby presentation, written as a
three-dimensional dotted--framed link
\[
  D_{\mathcal H}=D_{\mathrm{dot}}\sqcup D_{\mathrm{fr}}\subset S^3.
\]
An \emph{admissible section system} is a collection of mutually disjoint
cooriented spanning disks
\[
  \Delta_\delta\subset S^3,
  \qquad \partial\Delta_\delta=\delta,
\]
one for each dotted component, transverse to \(D_{\mathrm{fr}}\), together
with product neighborhoods of the disks.  Such systems exist because
\(D_{\mathrm{dot}}\) is an unlink.  Cutting the framed ribbon link along a
small product neighborhood
\(\Delta_\delta\times(-\epsilon,\epsilon)\) replaces every transverse
intersection by one endpoint on each of the two faces
\(\Delta_\delta^-\) and \(\Delta_\delta^+\).  The result is a
three-dimensional framed ribbon tangle with boundary.

Using the product coordinates, bunch the endpoints on each face into a
standard ordered row.  Their oriented strands form an ordered boundary word;
the opposite face carries its reversed orientation-dual word.  Regluing the
two faces by the identity matching of corresponding endpoints reconstructs
the original framed link.  No categorical label or morphism is part of this
topological construction.
\end{definition}

\begin{proposition}[Local movie normal forms]
\label{prop:local-movie-normal-forms}
The standard Kirby moves may be arranged as follows.
\begin{enumerate}[label=(\roman*),ref=\roman*]
  \item\label{item:slide-normal-form}
  After framed isotopy, every slide is supported in a \(3\)-ball and is a
  clean attachment, transport along the stationary component through
  finitely many local ribbon features, and a clean detachment.  For a
  \(1/1\) slide, the target dotted sublink is again an unlink and the source
  and target admit disjoint cross-sections agreeing outside the supporting
  ball.

  \item\label{item:cancellation-normal-form}
  A birth inserts an isolated standard cancelling pair.  Before a death,
  the remaining handles may be slid away from the pair so that the death
  removes an isolated standard pair.

  \item\label{item:fixed-exterior}
  Outside the supporting neighborhood, the framed tangle and all uninvolved
  components remain fixed.
\end{enumerate}
\end{proposition}

\begin{proof}
These are the standard band-slide and cancelling-pair models of Kirby
calculus; compare \cite[Sections~5.1 and~5.4]{GompfStipsicz1999} and
\cite[Table~2]{BarenzBarrett2018}.  Compact collars at the two ends of a
slide band can be made clean by a small isotopy.  Along the compact
stationary circle, generic position leaves only finitely many crossings,
extrema, and framing twists to pass.  For a cancelling \(1\)--\(2\) pair,
slide every other \(2\)-handle off the cancelling \(1\)-handle; the
\(2\)--\(3\) case is the upside-down analogue.  These reductions take place
inside the stated neighborhoods and leave the exterior fixed.
\end{proof}

\begin{lemma}[Cochain normal forms for cancelling pairs]
\label{lem:cancellation-cochain-normal-forms}
Let \(\mathcal H\to\mathcal H^+\) be a birth in the isolated standard local
form of \cref{prop:local-movie-normal-forms}.  For a $1$--$2$ birth,
extension by zero on the new primal $2$-handle sends a source $2$-cocycle
to a target $2$-cocycle representing the transported class.  Conversely,
every target $2$-cocycle is cohomologous to one whose new $2$-handle
coordinate is zero.  The same statement applies to an isolated $1$--$2$
death after the slide reduction in \cref{item:cancellation-normal-form}.
For a standard $2$--$3$ birth or death, every cocycle on the side containing
the pair has zero coordinate on the added $2$-handle.
\end{lemma}

\begin{proof}
After the other incidence entries have been cleared by the $2/2$ slides in
\cref{item:cancellation-normal-form}, the handle cochain complex of the
isolated cancelling $1$--$2$ pair contains the direct summand
\[
  \mathbb Z\xrightarrow{\,\pm1\,}\mathbb Z
\]
in degrees $1$ and $2$.  Extension by zero is the inclusion of the
complementary cochain complex.  Subtracting the coboundary of an appropriate
multiple of the new degree-$1$ basis cochain makes the new degree-$2$
coordinate zero without changing its cohomology class.

For a cancelling $2$--$3$ pair the corresponding direct summand lies in
degrees $2$ and $3$, again with differential $\pm1$.  The cocycle equation
in degree $3$ therefore forces the new degree-$2$ coordinate to vanish.
\end{proof}
